\subsection{Цели и прикладной смысл задачи синтеза новых видов}

\textbf{Задача синтеза новых видов (Novel View Synthesis)} заключается в создании изображений объекта или сцены с ракурсов, которые не были зафиксированы камерой изначально, на основе ограниченного набора существующих снимков.

Наиболее простым применением технологии является создание интерактивных туров или оцифровизация памятников \cite{11268436}. Например, можно отснять объект недвижимости, или памятник, и по видеоролику создать для конечного пользователя виртуальную его копию.

Её используют также в генерации данных для обучения. Например, в статье \cite{tancik2022blocknerfscalablelargescene} разработали способ, которых использует NeRF (один из подходов решения задачи NVS) для генерации цифровых двойников целых районов городов, которые можно использовать для обучения автопилотов. Сама специфика задачи позволяет генерировать фотореалистичные кадры, которые позволят автопилотам обучаться в приближенных к реальности условиях. При этом процесс построения модели NeRF автоматизирован.

NVS позволяет решать SLAM\footnote{SLAM \cite{slam} (англ. simultaneous localization and mapping — одновременная локализация и построение карты) — метод, используемый в мобильных автономных средствах для построения карты в неизвестном пространстве или для обновления карты в заранее известном пространстве с одновременным контролем текущего местоположения и пройденного пути.}, задачу локализации робота с одновременным построением карты местности, по данным с одной камеры (как например в статье \cite{matsuki2024gaussiansplattingslam}). Здесь уже сочетаются 2 особенности решения NVS, и 3DGS в частности: скорость, которая позволяет роботу картографировать местность практически в реальном времени, и дифференцируемость, которая позволяет сравнивать предполагаемые роботом координаты с тем, что видит камера, и таким образом корректировать накапливающиеся со временем расхождения.

В дополнение, так как решение NVS часто связано с созданием некоторого внутреннего представления сцены, как продукт из него можно попробовать извлекать карту глубины, нормалей, или даже 3D модель.